\section{Transferencia entre modalidades: primero representar el mundo como elementos que pueden comunicarse}

La familia de modelos responde «cómo pueden conectarse los token entre sí»; las aplicaciones entre modalidades dan un paso más y responden «qué cuenta como token». Karpathy resume la receta de transferencia con una frase deliberadamente exagerada: cortar imágenes, voz, trayectorias o estructuras moleculares en elementos, fingir que son texto y dejar que el Transformer modele sus relaciones. La verdadera dificultad de ingeniería se esconde justamente en el paso de «cortar en elementos».

\subsection{Vision Transformer: el patch es el token visual}

El capítulo anterior ya explicó que el encoder permite que un conjunto de token se comunique en ambas direcciones; el primer paso de la transferencia a visión es definir cómo convertir los píxeles bidimensionales en ese conjunto de token. ViT elige la interfaz más directa: divide la imagen en patch de tamaño fijo, aplana cada patch, lo proyecta linealmente a un vector y, tras agregar una representación de posición, lo envía al encoder. Si la imagen mide $H\times W$ y el lado del patch es $P$, el número de token es aproximadamente $N=HW/P^2$. Un patch más pequeño ofrece mayor granularidad, pero hace crecer rápidamente el cómputo $N^2$ de attention y la memoria de GPU; por eso la tokenización determina a la vez la resolución visual y el costo del sistema.

\lecturefigure{slide-47-vision-transformer.jpg}{Vision Transformer: la imagen se corta en una patch sequence y luego el encoder realiza la comunicación global.}{Video oficial de Stanford Online 00:52:41--00:53:42.}

\begin{knowledgebox}{Lectura de la figura: ViT no hace desaparecer la estructura bidimensional}
El patch ordering y el positional embedding siguen codificando la posición espacial, y el aumento de datos durante el entrenamiento y la escala de los datos también aportan bias visual. Un Transformer puro necesita volver a aprender la localidad y la regularidad ante traslaciones; por eso suele depender más que una ConvNet de los datos o del preentrenamiento. Sin embargo, la attention global permite que patch lejanos se comuniquen directamente desde las capas poco profundas y ofrece una interfaz unificada.
\end{knowledgebox}

\subsection{Speech: secuencias de spectrogram y bias convolucional local}

La voz puede convertirse primero en un spectrogram de time--frequency y luego cortarse a lo largo del tiempo en tramas o bloques. Conformer no consiste en «solo arrojar el espectrograma a un vanilla Transformer»: además de attention, agrega un convolution module para modelar al mismo tiempo las dependencias globales y los patrones acústicos locales; Whisper también usa una etapa de entrada especializada, encoder--decoder y un objetivo con supervisión débil a gran escala.

\lecturefigure{slide-48-conformer.jpg}{Conformer: combina el modelado global de attention con el bias acústico local de convolution.}{Video oficial de Stanford Online 00:53:42--00:54:00.}

\begin{knowledgebox}{Lectura de la figura: reutilizar entre dominios no significa que cada capa sea idéntica}
El espectrograma Mel es una representación de la energía de la voz en tiempo y frecuencia; la convolución local capta bien los patrones de tramas vecinas, y la autoatención capta bien las relaciones a larga distancia. La aportación de Conformer es precisamente combinar ambas. Lo que se dijo en clase, «tratarlo como texto», debe entenderse como compartir una interfaz de secuencia, no como borrar la estructura acústica ni las funciones objetivo especializadas.
\end{knowledgebox}

\subsection{Decision Transformer: escribir trayectorias de control como secuencias}

Decision Transformer intercala return-to-go, state y action, y deja que un causal model prediga la action condicionada al retorno esperado y al historial. Reescribe el offline RL como conditional sequence modeling: los datos de entrenamiento provienen de trayectorias existentes, sin interacción en línea con el entorno; la capacidad de planificación depende de la cobertura de los datos, del return conditioning y de que el modelo logre extraer una política de comportamiento a partir de las estadísticas de la secuencia.

\lecturefigure{slide-49-decision-transformer.jpg}{Decision Transformer: return, state y action forman una secuencia de trayectoria que puede modelarse de forma autorregresiva.}{Video oficial de Stanford Online 00:54:00--00:54:13.}

\begin{knowledgebox}{Lectura de la figura: no trata la recompensa como una palabra cualquiera}
Cada tipo de elemento tiene una semántica y un embedding distintos, y es necesario codificar el paso de tiempo y la información de modality/type. El modelo aprende una distribución condicional del tipo $p(a_t\mid R_t,s_{\le t},a_{<t})$, en lugar de resolver directamente la Bellman equation. El modelado de secuencias ofrece una interfaz de entrenamiento unificada, pero no elimina el offline data bias ni el riesgo de extrapolation.
\end{knowledgebox}

\subsection{AlphaFold2: attention dentro de un sistema especializado de estructura biológica}

El Evoformer de AlphaFold2 ejecuta varios tipos de attention y specialized updates entre la multiple-sequence alignment y la residue-pair representation; después, el módulo de estructura produce la conformación tridimensional. Esto muestra que attention puede manejar «qué residuos se relacionan con cuáles», pero no puede reducirse a un GPT que hace next-token prediction ordinaria sobre aminoácidos.

\lecturefigure{slide-50-alphafold.jpg}{AlphaFold2: attention se ubica dentro de un sistema de inferencia de estructura de proteínas altamente especializado en su dominio.}{Video oficial de Stanford Online 00:54:13--00:54:23.}

\begin{warningbox}{Error frecuente al leer la figura: confundir un backbone similar con sistemas equivalentes}
ViT, Conformer, Decision Transformer y AlphaFold2 usan attention, pero difieren en la unidad de entrada, la position geometry, la función objetivo, la combinación de módulos y la semántica de la salida. Lo que el curso quiere demostrar es que «la primitiva de comunicación direccionable es universal», no que «todos los problemas ya quedaron unificados por el modelado de lenguaje».
\end{warningbox}

\subsection{Entradas condicionales heterogéneas: por qué es fácil «agregar un sensor» a un Transformer}

El feature map de una red convolucional está estrechamente ligado a la geometría bidimensional; al agregar radar, map, vehicle state o audio hay que decidir en qué capa hacer concatenate o fuse. Un Transformer puede codificar cada tipo de entrada como elementos de un width común, agregar un type embedding para distinguir su origen y luego colocarlos en el conjunto donde se permite la comunicación; attention aprende pesos de fusión que dependen de los datos.

\lecturefigure{slide-51-tesla-input-flexibility.jpg}{Ejemplo de la experiencia en Tesla: la información heterogénea se codifica primero como elementos y luego se entrega a self-attention para fusionarla.}{Video oficial de Stanford Online 00:54:23--00:55:39.}

\begin{knowledgebox}{Términos clave: de dónde viene la libertad en las entradas}
\begin{itemize}
  \item \term{type embedding}: indica al modelo si el elemento proviene de image, radar, map o vehicle state;
  \item \term{positional encoding}: inyecta el tiempo, las coordenadas bidimensionales u otras relaciones geométricas;
  \item \term{mask/connectivity}: restringe qué elementos pueden comunicarse directamente;
  \item \term{tokenizer/adapter}: proyecta las distintas dimensiones originales a un residual width compartido;
  \item \term{inductive bias}: no desaparece, sino que se traslada del core block hacia la interfaz, las posiciones y las reglas de conexión.
\end{itemize}
\end{knowledgebox}

\teachervoice{El resumen oral de Karpathy es «cortar todo en pedazos, echarlo a la mezcla y dejar que self-attention decida cómo comunicarse». La frase capta la conveniencia de ingeniería, pero también se presta a malas interpretaciones. Una versión realmente reproducible debe responder: cómo cortar, cómo marcar el origen, qué aristas se permiten, con qué objetivo se supervisa y si el número de token es manejable.}

\subsection{Resumen de la sección}

El patrón común de la transferencia entre modalidades es representar los objetos del dominio como elementos con tipo y posición, y luego encaminar la información con attention. ViT sigue necesitando la posición espacial, Conformer conserva la convolución local, Decision Transformer depende del condicionamiento por trayectorias, AlphaFold2 usa representaciones especializadas de pair/MSA, y los sensores heterogéneos también requieren type embedding y connectivity. Un esqueleto compartido reduce la fricción de la fusión, pero no sustituye el modelado del dominio.
